\documentclass[11pt]{article}
\usepackage[margin=1in]{geometry}
\usepackage{amsmath,amssymb,microtype}
\usepackage[colorlinks=true,urlcolor=blue,citecolor=blue]{hyperref}
\title{The general-domain Taylor-centre conjecture of arXiv:1710.03114 was answered in 2019}
\author{Literature-status note for \texttt{fa\_banach\_001}}
\date{17 August 2026}

\begin{document}
\maketitle

\noindent\textbf{Status.}
\texttt{literature\_already\_answered (full conjecture, affirmative)}.

\section*{Original conjecture}

Let $\Omega\subset\mathbb C$ be a simply connected domain with
$0\notin\Omega$, and for $f\in\mathcal H(\Omega)$ put
\[
 T_Nf(z)=\sum_{n=0}^{N}\frac{f^{(n)}(z)}{n!}(-z)^n.
\]
The abstract and Section~1 of Panagiotis \cite{Panagiotis2019} record
Nestoridis's conjecture that
\[
 S(\Omega)=\{f\in\mathcal H(\Omega):\{T_Nf:N\geq0\}
                   \text{ is dense in }\mathcal H(\Omega)\}
\]
is a dense $G_\delta$ subset of $\mathcal H(\Omega)$.  That paper proves the
claim when $\Omega$ is a disc avoiding the origin, leaving the stated general
simply connected case open.

\section*{Exact later answer}

Bernal-Gonz\'alez, Cabana-M\'endez, Mu\~noz-Fern\'andez, and
Seoane-Sep\'ulveda \cite{BCMS2019} prove the full conjecture and more.
Their Theorem~2.10 (paper page~8; journal page~487) states that, for every
simply connected domain $G\subset\mathbb C$, the following are equivalent:
\begin{enumerate}
  \item $0\notin G$;
  \item the sequence $(T_N)$ is universal on $\mathcal H(G)$;
  \item $(T_N)$ is mixing;
  \item $S(G)$ is residual in $\mathcal H(G)$;
  \item $S(G)$ is dense-lineable in $\mathcal H(G)$.
\end{enumerate}
The same paper recalls that $S(G)$ is always a $G_\delta$ set; combined with
the residuality assertion, this gives precisely the dense $G_\delta$
conclusion requested in arXiv:1710.03114.  Thus the answer is
affirmative for every simply connected domain avoiding zero, with the stronger
mixing and dense-lineability conclusions.

The answering paper explicitly identifies Panagiotis's disc theorem in its
introduction and then announces that it will characterize non-vacuity for all
simply connected domains.  This is therefore a direct later answer, not an
indirect consequence inferred from unrelated machinery.

\section*{Independent route found during triage}

A direct attack before the final novelty check found a short alternative
coordinate picture.  For a branch $L=\log$ on $\Omega$, conjugation by
$f\mapsto f\circ\exp$ sends $T_N$ to
\[
 A_N=\prod_{j=1}^{N}\left(1-\frac{D}{j}\right),\qquad D=\frac{d}{dw}.
\]
On $e^{\lambda w}$, the eigenvalue is
$\prod_{j=1}^{N}(1-\lambda/j)$, which tends to zero for
$\operatorname{Re}\lambda>0$ and has modulus tending to infinity for
$\operatorname{Re}\lambda<0$.  Exponentials from either half-plane have
dense span in $\mathcal H(L(\Omega))$, giving the usual universality
criterion.  This corroborates the affirmative theorem but is not promoted as
a new result because Theorem~2.10 already settles the mathematical question.

\section*{Verification and scope}

The match was checked at the level of the operator formula, the domain
hypotheses, and the conclusion.  The original conjecture appears in the
abstract and Section~1 of arXiv:1710.03114.  The exact answer is
Theorem~2.10 of the 2019 paper.  Its Question~2.11 asks the strictly stronger
and separate question whether $S(G)$ is spaceable; no claim about that
question is made here.

\begin{thebibliography}{9}
\bibitem{Panagiotis2019}
C.~Panagiotis,
\emph{Universal partial sums of Taylor series as functions of the centre of
expansion}, Complex Variables and Elliptic Equations (2019),
\href{https://doi.org/10.1080/17476933.2019.1611791}{doi:10.1080/17476933.2019.1611791};
\href{https://arxiv.org/abs/1710.03114}{arXiv:1710.03114}.

\bibitem{BCMS2019}
L.~Bernal-Gonz\'alez, H.~J. Cabana-M\'endez, G.~A. Mu\~noz-Fern\'andez,
and J.~B. Seoane-Sep\'ulveda,
\emph{Universality of sequences of operators related to Taylor series},
J. Math. Anal. Appl. \textbf{474} (2019), 480--491,
\href{https://doi.org/10.1016/j.jmaa.2019.01.056}{doi:10.1016/j.jmaa.2019.01.056}.
\end{thebibliography}

\end{document}
